\documentclass[10pt]{article}
% Shared LaTeX style for MIT 18.06 exam revision notes.
% Compile topic files with Tectonic, XeLaTeX, or LuaLaTeX.

\usepackage[a4paper,margin=0.72in]{geometry}
\usepackage{fontspec}
\usepackage{amsmath,amssymb,mathtools}
\usepackage{booktabs,array,tabularx}
\usepackage{enumitem}
\usepackage{titlesec}
\usepackage{fancyhdr}
\usepackage{microtype}
\usepackage{xcolor}

\setmainfont{Times New Roman}
\setsansfont{Arial}

\setlength{\parindent}{0pt}
\setlength{\parskip}{3.6pt}
\setlength{\abovedisplayskip}{6pt}
\setlength{\belowdisplayskip}{6pt}
\setlength{\abovedisplayshortskip}{4pt}
\setlength{\belowdisplayshortskip}{4pt}
\renewcommand{\arraystretch}{1.08}
\setlength{\tabcolsep}{4.5pt}

\titleformat{\section}
  {\normalfont\large\bfseries}
  {\thesection}
  {0.55em}
  {}
\titlespacing*{\section}{0pt}{10pt}{3pt}

\titleformat{\subsection}
  {\normalfont\normalsize\bfseries}
  {\thesubsection}
  {0.5em}
  {}
\titlespacing*{\subsection}{0pt}{7pt}{2pt}

\setlist[itemize]{leftmargin=1.35em,itemsep=1pt,topsep=2pt,parsep=0pt}
\setlist[enumerate]{leftmargin=1.55em,itemsep=1pt,topsep=2pt,parsep=0pt}

\pagestyle{fancy}
\fancyhf{}
\fancyfoot[C]{\scriptsize\color{black!45}MIT 18.06 Linear Algebra Exam Notes --- \thepage}
\renewcommand{\headrulewidth}{0pt}
\renewcommand{\footrulewidth}{0pt}

\newcommand{\notesubtitle}[1]{%
  {\centering\itshape #1\par}
  \vspace{2pt}
}

\newcommand{\source}[1]{%
  {\centering\small #1\par}
  \vspace{6pt}
}

\newcommand{\labelnote}[2]{%
  \par\smallskip
  \noindent\textbf{#1.}\quad #2
  \par\smallskip
}

\newcommand{\tightlabel}[1]{%
  \par\smallskip
  \noindent\textbf{#1.}\quad
}

\newcolumntype{Y}{>{\raggedright\arraybackslash}X}
\newcolumntype{L}[1]{>{\raggedright\arraybackslash}p{#1}}

\newenvironment{notetable}
  {\small\begin{tabularx}{\linewidth}}
  {\end{tabularx}}


\begin{document}

\begin{center}
{\LARGE Complete Solutions Worked Example\par}
\end{center}
\notesubtitle{Concise revision notes for reading free variables, special solutions, and \(x=x_p+x_n\).}
\source{Based on handwritten MIT 18.06 revision notes.}

\section{Core Idea}

\labelnote{Core idea}{A system with free variables has infinitely many solutions. One particular solution solves \(Ax=b\); special solutions describe all movement inside \(N(A)\).}

The final answer has the form
\[
x=x_p+x_n.
\]
If \(s_1,\ldots,s_k\) are special solutions, then
\[
x=x_p+c_1s_1+\cdots+c_ks_k.
\]

\section{Reading Free Variables}

Free variables come from columns without pivots. In a reduced system, set each free variable one at a time to create nullspace directions.

\begin{center}
\small
\begin{tabularx}{\linewidth}{@{}L{0.28\linewidth}Y@{}}
\toprule
Variable type & Meaning \\
\midrule
Pivot variable & Forced by the equations. \\
Free variable & Chosen freely; creates a nullspace direction. \\
Particular choice & Usually set free variables to simple values to solve \(Ax=b\). \\
\bottomrule
\end{tabularx}
\end{center}

\section{Particular Solution}

For a consistent system
\[
Ax=b,
\]
choose convenient free variable values, often zeros, and solve for pivot variables. This gives one particular solution
\[
x_p.
\]
It is not the whole answer unless the nullspace is zero.

\section{Special Solutions}

To find the nullspace, solve
\[
Ax=0.
\]
For each free variable:
\begin{enumerate}
  \item set that free variable to \(1\);
  \item set the other free variables to \(0\);
  \item solve for pivot variables.
\end{enumerate}

The resulting vectors \(s_1,\ldots,s_k\) form a basis for \(N(A)\):
\[
N(A)=\operatorname{span}\{s_1,\ldots,s_k\}.
\]

\section{Putting the Answer Together}

If one particular solution is \(x_p\) and the nullspace basis is \(s_1,s_2\), write
\[
x=x_p+c_1s_1+c_2s_2.
\]
The constants \(c_1,c_2\) are arbitrary.

\labelnote{Meaning}{The particular solution reaches \(b\). The nullspace part moves without changing \(Ax\).}

\section{Common Mistakes}

\begin{center}
\small
\begin{tabularx}{\linewidth}{@{}L{0.40\linewidth}Y@{}}
\toprule
Mistake & Correct exam habit \\
\midrule
Giving only \(x_p\). & Add the full nullspace. \\
Using \(Ax=b\) to find special solutions. & Use \(Ax=0\). \\
Multiplying \(x_p\) by arbitrary constants. & Constants multiply only nullspace directions. \\
Forgetting one free variable. & Number of special solutions equals number of free variables. \\
Choosing complicated free values. & Use \(0\)'s and \(1\)'s unless the system suggests otherwise. \\
\bottomrule
\end{tabularx}
\end{center}

\section{Final Exam Checklist}

\tightlabel{Checklist}
\begin{enumerate}
  \item Row-reduce the augmented matrix \([A\ b]\).
  \item Check consistency.
  \item Mark pivot variables and free variables.
  \item Choose simple free-variable values to get \(x_p\).
  \item Solve \(Ax=0\) for special solutions.
  \item Create one special solution per free variable.
  \item Write \(x=x_p+c_1s_1+\cdots+c_ks_k\).
  \item Check by substituting: \(Ax_p=b\) and \(As_i=0\).
\end{enumerate}

\end{document}
